\chapter{Lung Surgery}

\section*{Definition and Overview}
Lung surgery encompasses a range of surgical procedures performed to diagnose, treat, or palliate diseases of the lungs and pleural space. The most common indication is the resection of primary lung malignancies, but surgery is also utilized for benign tumours, localized chronic infections, recurrent pneumothoraces, and diagnostic tissue biopsies. Surgical options range in extent from a wedge resection (removal of a small, non-anatomical wedge of lung tissue) and segmentectomy (removal of an anatomical segment) to lobectomy (removal of an entire lobe) and pneumonectomy (removal of an entire lung). Traditionally performed via an open thoracotomy, modern thoracic surgery increasingly utilises minimally invasive techniques, including Video-Assisted Thoracoscopic Surgery (VATS) and Robotic-Assisted Thoracoscopic Surgery (RATS), which significantly reduce post-operative pain, hospital stay, and recovery times.

\section*{Epidemiology}
Lung cancer is the leading cause of cancer-related mortality globally and in Australia, driving the majority of lung surgeries. Approximately 14,500 new cases of lung cancer are diagnosed in Australia annually, with surgery being the primary curative treatment for early-stage (Stage I to IIIA) non-small cell lung cancer (NSCLC). Pneumothorax, another common indication for surgical intervention (specifically VATS pleurodesis and bullectomy), has a bimodal distribution, frequently affecting tall, thin young males (primary spontaneous pneumothorax) or older individuals with underlying emphysema (secondary spontaneous pneumothorax). Surgical volume and patient outcomes are strongly associated with hospital volume and specialised thoracic surgical units.

\section*{Aetiology and Pathophysiology}
The indications for lung surgery arise from malignant transformations, structural defects, or chronic inflammatory processes in the pulmonary parenchyma and pleural cavity.
\begin{itemize}
    \item \textbf{Primary lung malignancies:} Squamous cell carcinoma, adenocarcinoma, and large cell carcinoma arising from the bronchial epithelium or alveolar cells, typically triggered by long-term exposure to tobacco smoke or occupational carcinogens.
    \item \textbf{Secondary metastases:} Metastatic spread of cancers from other primary sites (e.g. colorectal, breast, or renal cell carcinoma) to the lung parenchyma, where resection (metastasectomy) is indicated for oligometastatic disease.
    \item \textbf{Pleural bleb and bullae rupture:} Subpleural air-filled spaces that rupture into the pleural cavity, disrupting the negative intrapleural pressure and causing lung collapse (pneumothorax).
    \item \textbf{Chronic localized infections:} Destructive processes such as bronchiectasis, lung abscesses, or cavitary tuberculosis that fail to resolve with pharmacotherapy and require excision to prevent life-threatening haemoptysis or sepsis.
    \item \textbf{Severe emphysema:} Severe, heterogeneous tissue destruction leading to hyperinflation and mechanical disadvantage of the respiratory muscles, managed selectively with lung volume reduction surgery (LVRS).
\end{itemize}

\section*{Clinical Features}
The clinical features include pre-operative symptoms of the primary pathology and expected post-operative physical signs.
\begin{itemize}
    \item \textbf{Pre-operative symptoms:} Persistent cough, haemoptysis, dyspnoea, pleuritic chest pain, unexplained weight loss, and recurrent chest infections.
    \item \textbf{Surgical site pain:} Acute chest wall pain, which is typically severe after open thoracotomy due to rib spreading and intercostal nerve compression, and mild-to-moderate after VATS.
    \item \textbf{Altered respiratory parameters:} Reduced breath sounds on the operated side, mild tachypnoea, and transient oxygen desaturation requiring supplemental oxygen.
    \item \textbf{Chest tube drainage:} The presence of a chest drain emitting serosanguineous fluid and demonstrating a variable air leak (bubbling in the water-seal chamber).
    \item \textbf{Subcutaneous emphysema:} A crackling sensation (crepitus) felt under the skin of the chest wall or neck, caused by air tracking from the lung parenchyma into the subcutaneous tissues.
    \item \textbf{Coughing and sputum production:} Post-operative cough, occasionally productive of small amounts of blood-tinged sputum, as the remaining lung tissue expands and clears secretions.
\end{itemize}

\section*{Differential Diagnosis}
Post-operatively, clinicians must differentiate expected recovery findings from acute pulmonary and cardiovascular complications.
\begin{itemize}
    \item \textbf{Post-operative chest wall pain vs.\ myocardial infarction or pulmonary embolism:} Distinguishing surgical pain from ischaemic cardiac pain or the sudden pleuritic pain of a pulmonary embolism.
    \item \textbf{Atelectasis vs.\ pneumonia vs.\ pleural effusion:} Differentiating alveolar collapse (common in the first 48 hours) from active bacterial infection or fluid accumulation in the hemithorax.
    \item \textbf{Parenchymal air leak vs.\ chest tube connection leak:} Distinguishing a persistent leak from the sutured lung surface from a loose connection in the external drainage system.
    \item \textbf{Wound infection vs.\ normal post-operative erythema:} Differentiating standard inflammatory healing around incision/drain sites from cellulitis or deep wound infection.
\end{itemize}

\section*{When to Seek Medical Care}
Patients recovering from lung surgery should monitor their symptoms closely and seek medical review for worsening dyspnoea, fever, or increasing pain not managed by medications.

\textbf{Call 000 (or local emergency services) for:}
\begin{itemize}
    \item Sudden, severe shortness of breath or rapid, shallow breathing.
    \item New, sharp, stabbing chest pain, particularly if accompanied by sweating or lightheadedness.
    \item Coughing up large volumes of bright red blood (massive haemoptysis).
    \item Cyanosis (blue or grey discolouration of the lips, face, or fingernails).
    \item Sudden swelling of the neck or face accompanied by a crackling feeling under the skin and difficulty breathing.
    \item Dislodgement or blockage of the chest tube in a patient who has been discharged with a ambulatory drain.
\end{itemize}

\section*{Diagnosis and Investigations}
A thorough diagnostic and physiological workup is required to confirm the diagnosis and ensure the patient has sufficient cardiopulmonary reserve to survive the resection.
\begin{itemize}
    \item \textbf{Pulmonary function testing (PFTs):} Spirometry (measuring FEV1) and diffusing capacity of the lung for carbon monoxide (DLCO) to calculate the predicted post-operative lung function.
    \item \textbf{High-resolution CT scan of the chest:} Provides detailed anatomical mapping of the pulmonary nodules, lobar anatomy, and lymph node involvement.
    \item \textbf{Positron emission tomography (PET/CT):} Essential for staging lung cancer, evaluating nodal involvement, and screening for distant metastatic disease.
    \item \textbf{Cardiopulmonary exercise testing (CPET):} Performed in patients with borderline lung function to measure maximum oxygen consumption (VO2 max) and assess overall aerobic fitness.
    \item \textbf{Bronchoscopy with endobronchial ultrasound (EBUS):} Used to visualize the airway, obtain biopsy specimens of central lesions, and sample mediastinal lymph nodes for staging.
    \item \textbf{Arterial blood gas (ABG) analysis:} Evaluates baseline oxygenation, carbon dioxide levels, and acid-base status.
\end{itemize}

\section*{Management}
Management involves a combination of surgical intervention, post-operative chest tube care, aggressive pain control, and respiratory rehabilitation.
\begin{itemize}
    \item \textbf{Surgical resection:} Executing a lobectomy, segmentectomy, or wedge resection via thoracotomy, VATS, or RATS, ensuring clear oncological margins.
    \item \textbf{Chest tube management:} Insertion of one or two chest drains to evacuate air and blood, connected to a digital or water-seal suction system, which is monitored daily and removed once fluid output decreases and any air leak resolves.
    \item \textbf{Multimodal pain management:} Utilization of thoracic epidural analgesia, paravertebral nerve blocks, patient-controlled analgesia (PCA) with opioids, and scheduled oral paracetamol, NSAIDs, and neuropathic agents.
    \item \textbf{Aggressive respiratory physiotherapy:} Early mobilisation (out of bed on day 1), regular use of incentive spirometry, and assisted coughing techniques to prevent atelectasis and clear secretions.
    \item \textbf{Venous thromboembolism (VTE) prophylaxis:} Combination of sequential compression devices, early ambulation, and daily low-molecular-weight heparin (LMWH) injections.
    \item \textbf{Post-operative oncological follow-up:} Multi-disciplinary review to determine the need for adjuvant chemotherapy, immunotherapy, or radiotherapy.
\end{itemize}

\section*{Complications}
Lung surgery involves major physiological changes, exposing patients to serious post-operative complications.
\begin{itemize}
    \item \textbf{Prolonged air leak:} Air leaking from the lung surface into the pleural space for more than 5 to 7 days post-operatively, which is the most common minor complication.
    \item \textbf{Pneumonia and empyema:} Bacterial infection of the lung parenchyma or the pleural fluid, requiring targeted antibiotics and potentially repeat drainage.
    \item \textbf{Atrial fibrillation:} New-onset cardiac arrhythmia, highly common after thoracic surgery (especially pneumonectomy), managed with beta-blockers and anticoagulation.
    \item \textbf{Bronchopleural fistula:} A breakdown of the bronchial stump suture line, creating a communication between the airway and the pleural space, representing a life-threatening emergency.
    \item \textbf{Haemorrhage:} Intra-thoracic bleeding from major pulmonary vessels or intercostal arteries, which may necessitate urgent re-exploration.
    \item \textbf{Chylothorax:} Inadvertent injury to the thoracic duct leading to the accumulation of milky chyle in the pleural space, requiring dietary modifications or surgical ligation.
\end{itemize}

\section*{Prognosis and Outcomes}
The prognosis after lung surgery varies based on the underlying disease process, the extent of the resection, and the patient's pre-operative fitness. For early-stage (Stage IA) NSCLC, surgical resection offers an excellent 5-year survival rate of 80\% to 90\%. Minimally invasive VATS and RATS are associated with significantly lower perioperative morbidity, less chronic chest wall pain, and a faster return to baseline functional status compared to open thoracotomy. However, patients undergoing extensive resections (such as pneumonectomy) may experience a permanent reduction in exercise tolerance and require long-term cardiopulmonary adaptation.

\section*{Prevention and Self-Care}
Pre-operative preparation and dedicated post-operative self-care are essential to minimise complications and accelerate recovery.
\begin{itemize}
    \item \textbf{Smoking cessation:} Stopping smoking at least 4 weeks before surgery dramatically reduces the incidence of post-operative pulmonary complications and wound infections.
    \item \textbf{Incentive spirometry and deep breathing:} Performing deep breathing exercises 10 times every hour while awake to fully expand the remaining lung tissue.
    \item \textbf{Supported coughing:} Holding a pillow firmly against the chest incision (splinting) while coughing to reduce pain and support the chest wall.
    \item \textbf{Graduated physical activity:} Engaging in daily walking, avoiding heavy lifting (greater than 5 kg) for 4 to 6 weeks, and avoiding strenuous upper-body exercises until cleared by the surgeon.
    \item \textbf{Wound and drain care:} Keeping incisions clean and dry, monitoring for signs of infection, and managing ambulatory chest drains if discharged with one.
\end{itemize}

\section*{Sources and Further Reading}
The following organisations provide further information on lung surgery and thoracic health:
\begin{itemize}
    \item Cancer Council Australia. \textit{Understanding Lung Cancer and Surgical Options.} https://www.cancer.org.au/
    \item Thoracic Society of Australia and New Zealand (TSANZ). \textit{Resources for patients undergoing thoracic interventions.} https://www.thoracic.org.au/
    \item Mayo Clinic. \textit{Lobectomy and Lung Surgery: What to expect.} https://www.mayoclinic.org/
    \item NHS. \textit{Recovery after lung surgery.} https://www.nhs.uk/
    \item American Thoracic Society (ATS). \textit{Patient education: Thoracic surgery information.} https://www.thoracic.org/
\end{itemize}
